\subsection{Instrumenting A Function}

{

\#include <stdio.h>}





\bigskip



{

\#include {\textquotedbl}BPatch.h{\textquotedbl}}



{

\#include {\textquotedbl}BPatch_addressSpace.h{\textquotedbl}}



{

\#include {\textquotedbl}BPatch_process.h{\textquotedbl}}



{

\#include {\textquotedbl}BPatch_binaryEdit.h{\textquotedbl}}



{

\#include {\textquotedbl}BPatch_point.h{\textquotedbl}}



{

\#include {\textquotedbl}BPatch_function.h{\textquotedbl}}





\bigskip



{

using namespace std;}



{

using namespace Dyninst;}





\bigskip



{

// Create an instance of class BPatch}



{

BPatch bpatch;}





\bigskip



{

// Different ways to perform instrumentation}



{

typedef enum {}



{

create,}



{

attach,}



{

open}



{

} accessType_t;





\bigskip



{

// Attach, create, or open a file for rewriting}



{

BPatch_addressSpace* startInstrumenting(accessType_t accessType,}



{

const char* name,}



{

int pid,}



{

const char* argv[]) {}



{

BPatch_addressSpace* handle = NULL;}





\bigskip



{

switch(accessType) {}



{

case create:}



{

handle = bpatch.processCreate(name, argv);}



{

if (!handle) { fprintf(stderr, {\textquotedbl}processCreate

failed{\textbackslash}n{\textquotedbl}); }}



{

break;}



{

case attach:}



{

handle = bpatch.processAttach(name, pid);}



{

if (!handle) { fprintf(stderr, {\textquotedbl}processAttach

failed{\textbackslash}n{\textquotedbl}); }}



{

break;}



{

case open:}



{

// Open the binary file and all dependencies}



{

handle = bpatch.openBinary(name, true);}



{

if (!handle) { fprintf(stderr, {\textquotedbl}openBinary

failed{\textbackslash}n{\textquotedbl}); }}



{

break;}



{

}}





\bigskip



{

return handle;}



{

}}





\bigskip





\bigskip





\bigskip





\bigskip



{

// Find a point at which to insert instrumentation}



{

std::vector<BPatch_point*>* findPoint(BPatch_addressSpace* app,



{

const char* name,



{

BPatch_procedureLocation loc) {}



{

std::vector<BPatch_function*> functions;}



{

std::vector<BPatch_point*>* points;}





\bigskip



{

// Scan for functions named {\textquotedbl}name{\textquotedbl}}



{

BPatch_image* appImage = app->getImage();}



{

appImage->findFunction(name, functions);}



{

if (functions.size() == 0) {}



{

fprintf(stderr, {\textquotedbl}No function \%s{\textbackslash}n{\textquotedbl}, name);}



{

return points;}



{

} else if (functions.size() > 1) {}



{

fprintf(stderr, {\textquotedbl}More than one \%s; using the first one{\textbackslash}n{\textquotedbl},

name);}



{

}}





\bigskip



{

// Locate the relevant points}



{

points = functions[0]->findPoint(loc);}



{

return points;}



{

}}





\bigskip



{

// Create and insert an increment snippet}



{

bool createAndInsertSnippet(BPatch_addressSpace* app,}



{

std::vector<BPatch_point*>* points) {}



{

BPatch_image* appImage = app->getImage();}





\bigskip



{

// Create an increment snippet}



{

BPatch_variableExpr* intCounter =



{

app->malloc(*(appImage->findType({\textquotedbl}int{\textquotedbl})),

{\textquotedbl}myCounter{\textquotedbl});}



{

BPatch_arithExpr addOne(BPatch_assign,}



{

*intCounter,}



{

BPatch_arithExpr(BPatch_plus,



{

*intCounter,



{

BPatch_constExpr(1)));}





\bigskip



{

// Insert the snippet}



{

if (!app->insertSnippet(addOne, *points)) {}



{

fprintf(stderr, {\textquotedbl}insertSnippet failed{\textbackslash}n{\textquotedbl});}



{

return false;}



{

}}



{

return true;}



{

}}





\bigskip



{

// Create and insert a printf snippet}



{

bool createAndInsertSnippet2(BPatch_addressSpace* app,}



{

std::vector<BPatch_point*>* points) {}



{

BPatch_image* appImage = app->getImage();}





\bigskip



{

// Create the printf function call snippet}



{

std::vector<BPatch_snippet*> printfArgs;}



{

BPatch_snippet* fmt =



{

new BPatch_constExpr({\textquotedbl}InterestingProcedure called \%d

times{\textbackslash}n{\textquotedbl});}



{

printfArgs.push_back(fmt);}





\bigskip



{

BPatch_variableExpr* var = appImage->findVariable({\textquotedbl}myCounter{\textquotedbl});}



{

if (!var) {}



{

fprintf(stderr, {\textquotedbl}Could not find 'myCounter' variable{\textbackslash}n{\textquotedbl});}



{

return false;}



{

} else {}



{

printfArgs.push_back(var);}



{

}}





\bigskip



{

// Find the printf function}



{

std::vector<BPatch_function*> printfFuncs;}



{

appImage->findFunction({\textquotedbl}printf{\textquotedbl}, printfFuncs);}



{

if (printfFuncs.size() == 0) {}



{

fprintf(stderr, {\textquotedbl}Could not find printf{\textbackslash}n{\textquotedbl});}



{

return false;}



{

}}





\bigskip



{

// Construct a function call snippet}



{

BPatch_funcCallExpr printfCall(*(printfFuncs[0]), printfArgs);}





\bigskip



{

// Insert the snippet}



{

if (!app->insertSnippet(printfCall, *points)) {}



{

fprintf(stderr, {\textquotedbl}insertSnippet failed{\textbackslash}n{\textquotedbl});}



{

return false;}



{

}}



{

return true;}



{

}}





\bigskip



{

void finishInstrumenting(BPatch_addressSpace* app, const char* newName)}



{

{}



{

BPatch_process* appProc = dynamic_cast<BPatch_process*>(app);}



{

BPatch_binaryEdit* appBin = dynamic_cast<BPatch_binaryEdit*>(app);}





\bigskip



{

if (appProc) {}



{

if (!appProc->continueExecution()) {}



{

fprintf(stderr, {\textquotedbl}continueExecution failed{\textbackslash}n{\textquotedbl});}



{

}}



{

while (!appProc->isTerminated()) {}



{

bpatch.waitForStatusChange();}



{

}}



{

} else if (appBin) {}



{

if (!appBin->writeFile(newName)) {}



{

fprintf(stderr, {\textquotedbl}writeFile failed{\textbackslash}n{\textquotedbl});}



{

}}



{

}}



{

}}





\bigskip



{

int main() {}



{

// Set up information about the program to be instrumented}



{

const char* progName = {\textquotedbl}InterestingProgram{\textquotedbl};}



{

int progPID = 42;}



{

const char* progArgv[] = {{\textquotedbl}InterestingProgram{\textquotedbl}, {\textquotedbl}-h{\textquotedbl},

NULL};}



{

accessType_t mode = create;}





\bigskip



{

// Create/attach/open a binary}



{

BPatch_addressSpace* app =



{

startInstrumenting(mode, progName, progPID, progArgv);}



{

if (!app) {}



{

fprintf(stderr, {\textquotedbl}startInstrumenting failed{\textbackslash}n{\textquotedbl});}



{

exit(1);}



{

}}





\bigskip



{

// Find the entry point for function InterestingProcedure}



{

const char* interestingFuncName = {\textquotedbl}InterestingProcedure{\textquotedbl};}



{

std::vector<BPatch_point*>* entryPoint =



{

findPoint(app, interestingFuncName, BPatch_entry);}



{

if (!entryPoint {\textbar}{\textbar} entryPoint->size() == 0) {}



{

fprintf(stderr, {\textquotedbl}No entry points for \%s{\textbackslash}n{\textquotedbl},

interestingFuncName);}



{

exit(1);}



{

}}





\bigskip



{

// Create and insert instrumentation snippet}



{

if (!createAndInsertSnippet(app, entryPoint)) {}



{

fprintf(stderr, {\textquotedbl}createAndInsertSnippet failed{\textbackslash}n{\textquotedbl});}



{

exit(1);}



{

}}





\bigskip



{

// Find the exit point of main}



{

std::vector<BPatch_point*>* exitPoint =



{

findPoint(app, {\textquotedbl}main{\textquotedbl}, BPatch_exit);}



{

if (!exitPoint {\textbar}{\textbar} exitPoint->size() == 0) {}



{

fprintf(stderr, {\textquotedbl}No exit points for main{\textbackslash}n{\textquotedbl});}



{

exit(1);}



{

}}





\bigskip



{

// Create and insert instrumentation snippet 2}



{

if (!createAndInsertSnippet2(app, exitPoint)) {}



{

fprintf(stderr, {\textquotedbl}createAndInsertSnippet2 failed{\textbackslash}n{\textquotedbl});}



{

exit(1);}



{

}}



{

// Finish instrumentation}



{

const char* progName2 = {\textquotedbl}InterestingProgram-rewritten{\textquotedbl};}



{

finishInstrumenting(app, progName2);}



{

}}



